% GENERATED FILE. Edit canonical lessons, then run npm run generate:books.
% canonical-source-hash: fnv1a64:ad80a2a9ba5d2375
% canonical-lessons: RU-C520-rezhissyor, RU-C520-aktrisa, RU-C520-rol, RU-C520-zritel, RU-C520-snimatsya

\chapter{Directors and Actresses}
\label{ch:ru-directors-and-actresses}
\hlchaptermodality{520}

\begin{chapteropening}
\textbf{By the end of this chapter:} \emph{I can talk about the people who make a film.}

Complete the last lesson of chapter 520: I can talk about the people who make a film.
\end{chapteropening}

\section[\texorpdfstring{\ru{режиссёр} (rezhissyór)}{режиссёр (rezhissyór)}]{\ru{режиссёр} (rezhissyór) — a director (of a film or play)}

\label{lesson:RU-C520-rezhissyor}



\begin{quote}
Before the new one: say the Russian for a premiere, then the Russian for to bow.
\end{quote}

\subsection*{You'll want to know: \ru{режиссёр}}
\textbf{\ru{режиссёр}} (\emph{rezhissyór}) — \textquotedblleft{}a director (of a film or play)\textquotedblright{}.

\begin{quote}
\ru{Режиссёр} \ru{долго} \ru{выбирал} \ru{актёров}.
\end{quote}

\begin{quote}
\emph{Rezhissyór dólgo vybirál aktyórov.}
\end{quote}

\begin{quote}
The director took a long time choosing the actors.
\end{quote}

Masculine; also of a woman.

\begin{cousinweb}
\emph{\ru{режиссёр}} was borrowed from French \emph{régisseur} 'a manager, a stage manager' (EtymDB).
\end{cousinweb}

\begin{grammarlens}[title={What you've built}]
One more word for telling, explaining and describing at B1.
\end{grammarlens}

\subsection*{Your turn}
\begin{itemize}
\raggedright
  \item \emph{Say it:} \emph{rezhissyór}
  \item \emph{Say it:} \emph{rezhissyór}, once more
  \item \emph{Say it:} a sentence of your own with \emph{rezhissyór}
  \item \emph{Recall:} say the Russian for a premiere, then the Russian for to bow, then say \emph{rezhissyór} again
\end{itemize}

\begin{tcolorbox}[breakable,colback=teal!4,colframe=teal!35!black,title={Before you move on}]
What does \ru{режиссёр} mean? (\textquotedblleft{}A director (of a film or play)\textquotedblright{}.) Say it once more.
\end{tcolorbox}
\section[\texorpdfstring{\ru{актриса} (aktrísa)}{актриса (aktrísa)}]{\ru{актриса} (aktrísa) — an actress}

\label{lesson:RU-C520-aktrisa}



\begin{quote}
Before the new one: say the Russian for to bow, then the Russian for a director.
\end{quote}

\subsection*{You'll want to know: \ru{актриса}}
\textbf{\ru{актриса}} (\emph{aktrísa}) — \textquotedblleft{}an actress\textquotedblright{}.

\begin{quote}
\ru{Эта} \ru{актриса} \ru{часто} \ru{играет} \ru{в} \ru{комедиях}.
\end{quote}

\begin{quote}
\emph{Éta aktrísa chásto igráyet v komédiyakh.}
\end{quote}

\begin{quote}
This actress often appears in comedies.
\end{quote}

Feminine.

\begin{cousinweb}
\emph{\ru{актриса}} is the feminine of \emph{\ru{актёр}} 'actor', usually traced to French \emph{actrice}, ultimately from Latin \emph{āctor}; EtymDB lists \emph{\ru{актриса}} without a source.
\end{cousinweb}

\begin{grammarlens}[title={What you've built}]
One more word for telling, explaining and describing at B1.
\end{grammarlens}

\subsection*{Your turn}
\begin{itemize}
\raggedright
  \item \emph{Say it:} \emph{aktrísa}
  \item \emph{Say it:} \emph{aktrísa}, once more
  \item \emph{Say it:} a sentence of your own with \emph{aktrísa}
  \item \emph{Recall:} say the Russian for to bow, then the Russian for a director, then say \emph{aktrísa} again
\end{itemize}

\begin{tcolorbox}[breakable,colback=teal!4,colframe=teal!35!black,title={Before you move on}]
What does \ru{актриса} mean? (\textquotedblleft{}An actress\textquotedblright{}.) Say it once more.
\end{tcolorbox}
\section[\texorpdfstring{\ru{роль} (rol)}{роль (rol)}]{\ru{роль} (rol) — a role, a part}

\label{lesson:RU-C520-rol}



\begin{quote}
Before the new one: say the Russian for a director, then the Russian for an actress.
\end{quote}

\subsection*{You'll want to know: \ru{роль}}
\textbf{\ru{роль}} (\emph{rol}) — \textquotedblleft{}a role, a part\textquotedblright{}.

\begin{quote}
\ru{Он} \ru{играет} \ru{главную} \ru{роль} \ru{в} \ru{новом} \ru{фильме}.
\end{quote}

\begin{quote}
\emph{On igráyet glávnuyu rol v nóvom fílme.}
\end{quote}

\begin{quote}
He plays the lead role in the new film.
\end{quote}

Feminine; genitive \ru{роли}.

\begin{cousinweb}
\emph{\ru{роль}} was borrowed from French \emph{rôle} and German \emph{Rolle}, which go back to Latin \emph{rotulus} 'a small roll, a scroll' (EtymDB), the actor's roll of lines.
\end{cousinweb}

\begin{grammarlens}[title={What you've built}]
One more word for telling, explaining and describing at B1.
\end{grammarlens}

\subsection*{Your turn}
\begin{itemize}
\raggedright
  \item \emph{Say it:} \emph{rol}
  \item \emph{Say it:} \emph{rol}, once more
  \item \emph{Say it:} a sentence of your own with \emph{rol}
  \item \emph{Recall:} say the Russian for a director, then the Russian for an actress, then say \emph{rol} again
\end{itemize}

\begin{tcolorbox}[breakable,colback=teal!4,colframe=teal!35!black,title={Before you move on}]
What does \ru{роль} mean? (\textquotedblleft{}A role, a part\textquotedblright{}.) Say it once more.
\end{tcolorbox}
\section[\texorpdfstring{\ru{зритель} (zrítel)}{зритель (zrítel)}]{\ru{зритель} (zrítel) — a viewer, a spectator}

\label{lesson:RU-C520-zritel}



\begin{quote}
Before the new one: say the Russian for an actress, then the Russian for a role.
\end{quote}

\subsection*{You'll want to know: \ru{зритель}}
\textbf{\ru{зритель}} (\emph{zrítel}) — \textquotedblleft{}a viewer, a spectator\textquotedblright{}.

\begin{quote}
\ru{Фильм} \ru{понравился} \ru{зрителям}.
\end{quote}

\begin{quote}
\emph{Film ponrávilsya zrítelyam.}
\end{quote}

\begin{quote}
Viewers liked the film.
\end{quote}

Masculine; feminine \ru{зрительница}. \emph{\ru{Зрительный} \ru{зал}} — the auditorium.

\begin{cousinweb}
\emph{\ru{зритель}} is built on \emph{\ru{зреть}} 'to see' with \emph{-\ru{тель}}; EtymDB lists \emph{\ru{зритель}} without a source.
\end{cousinweb}

\begin{grammarlens}[title={What you've built}]
One more word for telling, explaining and describing at B1.
\end{grammarlens}

\subsection*{Your turn}
\begin{itemize}
\raggedright
  \item \emph{Say it:} \emph{zrítel}
  \item \emph{Say it:} \emph{zrítel}, once more
  \item \emph{Say it:} a sentence of your own with \emph{zrítel}
  \item \emph{Recall:} say the Russian for an actress, then the Russian for a role, then say \emph{zrítel} again
\end{itemize}

\begin{tcolorbox}[breakable,colback=teal!4,colframe=teal!35!black,title={Before you move on}]
What does \ru{зритель} mean? (\textquotedblleft{}A viewer, a spectator\textquotedblright{}.) Say it once more.
\end{tcolorbox}
\section[\texorpdfstring{\ru{сниматься} (snimátsya)}{сниматься (snimátsya)}]{\ru{сниматься} (snimátsya) — to act in a film; to have one's photo taken}

\label{lesson:RU-C520-snimatsya}



\begin{quote}
Before the new one: say the Russian for a role, then the Russian for a viewer.
\end{quote}

\subsection*{You'll want to know: \ru{сниматься}}
\textbf{\ru{сниматься}} (\emph{snimátsya}) — \textquotedblleft{}to act in a film; to have one's photo taken\textquotedblright{}.

\begin{quote}
\ru{Она} \ru{снимается} \ru{в} \ru{кино} \ru{с} \ru{детства}.
\end{quote}

\begin{quote}
\emph{Oná snimáyetsya v kinó s détstva.}
\end{quote}

\begin{quote}
She has been acting in films since childhood.
\end{quote}

Imperfective; perfective \ru{сняться} (\ru{снимусь}). \emph{\ru{Сниматься} \ru{в} \ru{фильме}} — to be in a film.

\begin{cousinweb}
\emph{\ru{сниматься}} is the reflexive of \emph{\ru{снимать}} 'to take off; to film', built on \emph{\ru{с}-} and the root \emph{-\ru{ним}-/-\ru{ня}-} 'to take'; EtymDB has no entry for the reflexive.
\end{cousinweb}

\begin{grammarlens}[title={What you've built}]
One more word for telling, explaining and describing at B1.
\end{grammarlens}

\subsection*{Your turn}
\begin{itemize}
\raggedright
  \item \emph{Say it:} \emph{snimátsya}
  \item \emph{Say it:} \emph{snimátsya}, once more
  \item \emph{Say it:} a sentence of your own with \emph{snimátsya}
  \item \emph{Recall:} say the Russian for a role, then the Russian for a viewer, then say \emph{snimátsya} again
\end{itemize}

\begin{tcolorbox}[breakable,colback=teal!4,colframe=teal!35!black,title={Before you move on}]
What does \ru{сниматься} mean? (\textquotedblleft{}To act in a film; to have one's photo taken\textquotedblright{}.) Say it once more.
\end{tcolorbox}
